\chapter{Metodyka pracy}
\label{cha:methodology}

Rozdział ten poświęcony jest metodyce pracy przyjętej podczas realizacji niniejszej pracy dyplomowej. Znajdują się w nim informacje o narzędziach wspomagających proces twórczy, ze szczególnym uwzględnieniem wykorzystania sztucznej inteligencji, a także wyszczególnienie samodzielnego wkładu autora. Ma to na celu transparentne przedstawienie sposobu, w jaki powstał opisany system oraz rzetelne oddzielenie zautomatyzowanej pomocy od pracy własnej dyplomanta.

\section{Wykorzystanie wsparcia sztucznej inteligencji}

Złożoność oraz rozległy zakres technologiczny projektu - obejmujący zarówno tworzenie interfejsu użytkownika, jak i logiki po stronie serwera - skłoniły dyplomanta do zastosowania nowoczesnych narzędzi opartych na sztucznej inteligencji w procesie programistycznym. Wsparcie to obejmowało korzystanie z tradycyjnych asystentów AI, jak również autonomicznych agentów AI wykonujących powierzone sekwencje zadań technicznych (takich jak analiza kodu w wielu plikach nawzajem). Podczas tworzenia wykorzystywane były w szczególności modele z rodziny Gemini (wersje 3 oraz 3.1 Pro).

Zastosowanie tych technologii pozwoliło na przyspieszenie powtarzalnych operacji i szybsze odnajdywanie specyficznych błędów w kodzie, w związku z czym modele te pełniły funkcję wirtualnego asystenta. Głównymi obszarami, w których posłużono się generatywną sztuczną inteligencją, były:
\begin{itemize}
    \item Generowanie kodu pomocniczego \cite{boilerplate_code}: Zautomatyzowanie tworzenia powtarzalnych i nużących struktur, deklaracji podstawowych widoków frontendowych czy bazowych szablonów dla plików konfiguracyjnych.
    \item Debugowanie: Skonsultowanie trudnych do zinterpretowania logów, ostrzeżeń z bibliotek czy błędów pojawiających się w konsoli, co znacząco zmniejszyło czas potrzebny na lokalizację usterek.
    \item Konsultacje technologiczne i dokumentacyjne: Szybkie odszukiwanie właściwych metod w nowych bibliotekach czy przypominanie rzadko używanych fragmentów składni.
\end{itemize}

\section{Wkład autorski i decyzyjność}

Mimo wsparcia ze strony narzędzi AI, należy podkreślić, że stanowiły one wyłącznie rozwiązanie wspomagające. Cały system w swoim ostatecznym kształcie jest wynikiem samodzielnej pracy koncepcyjnej, badawczej i programistycznej dyplomanta.

Do kluczowych zadań inżynierskich, zrealizowanych całkowicie samodzielnie przez autora, należą:
\begin{itemize}
    \item Projektowanie architektury: Samodzielne zaprojektowanie podziału na logiczne warstwy usług, wybór mechanizmów komunikacji między mikroserwisami (\gls{REST} \gls{API}, WebSocket) oraz decyzyjność co do stosu technologicznego (FastAPI, Redis, Celery, React, baza PostgreSQL).
    \item Implementacja silnika generowania danych: Zaprojektowanie rdzenia logiki backendowej odpowiadającej za walidację oraz generację syntetycznych danych. Wyróżnia się tu m.in. implementację algorytmu sortowania topologicznego rozwiązującego zależności pomiędzy rozbudowanymi relacyjnie tabelami baz danych.
    \item Projekt interfejsu i interakcji użytkownika (\gls{UI}/\gls{UX}): Opracowanie unikalnej koncepcji użytej aplikacji webowej i wizualizacji układu paneli roboczych, która musiała zapewniać ergonomię obsługi.
    \item Integracja systemu i ocena kodu: Każdy fragment kodu zaproponowany przez narzędzia AI weryfikowany był przez autora pod kątem poprawności i zgodności z ogólną koncepcją systemu.
\end{itemize}

Wykorzystanie nowoczesnych asystentów programistycznych znacząco odciążyło dyplomanta z najbardziej mechanicznych oraz trywialnych aspektów pisania kodu oprogramowania. Zaoszczędzony w ten sposób czas pozwolił na znacznie głębsze skupienie się na meritum pracy: rozwiązywaniu problemów architektonicznych, a także dopracowywaniu detali implementacyjnych.
